\documentclass{imc-inf}

\title{Autonomous Perception-Fused Driving on the Jetson Orin Nano Platform}
\subtitle{Camera Fusion for Perception and Control --- Waveshare JETANK AI Kit + NVIDIA Jetson Orin Nano Super}
\thesistype{Capstone Project Report}
\author{Adam Ryan, Aidan Ryan, Denis Vasilev, Nora Regina Cerveny}
\supervisor{Prof. Tingting Yuan}
\copyrightyear{2026}
\submissiondate{26.06.2026}
\keywords{Jetson Orin Nano, JETANK, Autonomous Driving, Camera Fusion, Edge AI, Robotics}

\usepackage{booktabs}
\usepackage{longtable}
\usepackage{listings}
\usepackage{array}
\usepackage{textcomp}

\lstset{
    basicstyle=\footnotesize\ttfamily,
    frame=single,
    breaklines=true,
    captionpos=b,
    tabsize=2
}

\begin{document}
\frontmatter\maketitle{}

%% --- CUSTOM DECLARATION (group version) -----------------------------------
\chapter*{Declaration of honour}
\addtoToC{Declaration of honour}

We declare on our word of honour that we have written this \thesistype{} on our own and that we have not used any sources or resources other than stated and that we have marked those passages and/or ideas that were either verbally or textually extracted from sources. This also applies to drawings, sketches, graphic representations as well as to sources from the internet. The \thesistype{} has not been submitted in this or similar form for assessment at any other domestic or foreign post-secondary educational institution and has not been published elsewhere. The present \thesistype{} complies with the version submitted electronically.

\vspace{3em}

\begin{tabular}{@{}p{.5in}p{3in}@{}}
 & \hrulefill \\
 & Adam Ryan \\
 & \\[1.5em]
 & \hrulefill \\
 & Aidan Ryan \\
 & \\[1.5em]
 & \hrulefill \\
 & Denis Vasilev \\
 & \\[1.5em]
 & \hrulefill \\
 & Nora Regina Cerveny \\
\end{tabular}

\newpage
%% ---------------------------------------------------------------------------

\begin{abstract}
This report documents a semester-long capstone project to build an autonomous,
camera-fusion-driven mobile robot on a Waveshare JETANK AI Kit chassis retrofitted
with an NVIDIA Jetson Orin Nano Super compute module. The project was scoped into
three tasks: Task~C (hardware assembly and platform integration), Task~B (camera
fusion for perception, combining lane detection and object detection into a single
per-frame state), and Task~A (camera fusion for driving, converting that fused state
into safety-gated steering and throttle commands).

The team's core technical finding is that the JETANK chassis, designed around the
original Jetson Nano, is not a drop-in match for the Jetson Orin Nano Super. The
two platforms differ in power delivery, kernel packaging, camera driver model,
I2C bus numbering, network interface naming, and motor-driver channel mapping.

Across the semester, ten distinct system-level integration issues were identified
and resolved, spanning OS imaging, WiFi driver support, desktop environment, Python
packaging, CSI camera enablement, OLED display, and serial servo calibration. Two
blockers were not resolved before the project deadline: a power architecture mismatch
requiring an independently sourced power bank, and unresponsive drive motors where
two confirmed software bugs were fixed but the root cause remains at the hardware
layer. As a direct consequence, Tasks~B and~A were fully designed at the architectural
level but could not be implemented or tested on hardware.

This report presents the working hypothesis for the remaining motor fault,
a documented bring-up path for Orin-generation JetBot hardware, and the complete
design for the fused-perception and driving-control pipelines.
\end{abstract}



\addtoToC{Table of Contents}%
\tableofcontents%
\clearpage

\addtoToC{List of Tables}%
\listoftables
\clearpage



%%   MAIN MATTER  %%%%%%%%%%%%%%%%%%%%%%%%%%%%%%%%%%%%%%%%%%%%%%%%%%%%%%%%%%%%%%
\mainmatter%

%==============================================================================
\chapter{Introduction}
\label{chap:introduction}
%==============================================================================

\section{Why Edge AI Robotics}
\label{sec:why-edge-ai}

Autonomous mobile robots that perform inference locally, rather than relying on a
cloud connection, are essential wherever latency, privacy, or connectivity cannot
be guaranteed. Warehouses, hospitals, agricultural inspection, and last-mile
delivery are common use cases. The NVIDIA Jetson family of single-board computers
provides GPU-accelerated edge inference in a form factor compact enough for small
robots, making it a natural teaching platform for this kind of system \cite{nvidia_jetbot,nvidia_jetson_orin}.

\section{Why Camera Fusion}
\label{sec:why-camera-fusion}

Individual perception modalities are fragile in isolation. Lane detection alone
fails the moment an obstacle obscures the road markings; object detection alone
cannot determine where the robot sits within its lane. Fusing both into a single
per-frame state (lane offset plus obstacle class, position, and distance)
produces a more complete and robust basis for control decisions than either signal
alone. This is the central technical idea the project set out to demonstrate.

\section{The JetBot / Orin Generation Gap}
\label{sec:generation-gap}

The widely used NVIDIA JetBot reference platform targets the original Jetson Nano
(2019, JetPack~4.x). This project instead used the newer Jetson Orin Nano Super
(2024, JetPack~6.x), mounted on a Waveshare JETANK chassis that was itself designed
around the older Nano \cite{waveshare_jetank}. The two Jetson generations differ substantially in power
requirements, kernel packaging, camera driver model, and peripheral addressing.
Existing JetBot and JETANK tutorials assume the older hardware and require substantial
adaptation at nearly every layer of the stack. A secondary contribution of this project,
beyond the three assigned tasks, is a documented bring-up path for teams attempting
the same Orin-generation upgrade in the future.

\section{Task Scope}
\label{sec:task-scope}

\begin{description}
\item[Task C --- Hardware Assembly and Platform Integration]
Chassis build, compute module mounting, camera/motor/display validation, and
remote development workflow.
\item[Task B --- Camera Fusion for Perception]
Combine a lane-following model and a YOLO-based object detector \cite{redmon2016yolo,ultralytics_yolo} into one
structured, per-frame perception state.
\item[Task A --- Camera Fusion for Driving]
Convert the fused perception state into a safety-gated control policy
(steering correction, speed control, emergency stop).
\end{description}

%==============================================================================
\chapter{Materials and System Overview}
\label{chap:materials}
%==============================================================================

This chapter lists the hardware and software configuration as of the end of the
semester, including the corrected channel mappings and bus numbers discovered
during bring-up (Chapter~\ref{chap:task-c}).

\begin{table}[htbp]
\centering
\caption{Hardware components used in the project.}
\label{tab:hardware}
\footnotesize
\begin{tabular}{@{}p{0.9in}p{2.6in}p{1.8in}@{}}
\toprule
\textbf{Component} & \textbf{Model / Spec} & \textbf{Purpose} \\
\midrule
Main compute & NVIDIA Jetson Orin Nano Super (8~GB) & GPU-accelerated inference;
6$\times$ ARM Cortex-A78, Ampere GPU, 1024 CUDA cores \\
Chassis & Waveshare JETANK AI Kit & Tank-track chassis; TB6612FNG dual
H-bridge motor driver \cite{toshiba_tb6612fng}; PCA9685 PWM controller \cite{nxp_pca9685} \\
Camera & Waveshare IMX219-160 (CSI) & 160$^\circ$ FOV, 8~MP, MIPI CSI-2 to CAM0 \\
Servos & Feetech SCS15-AP $\times$5 & Serial bus servos for 5-DOF arm
(2 mounted: base rotation + camera tilt) \\
Display & SSD1306 OLED (128$\times$32) & System status: IP address, CPU/RAM,
component health \\
Power (Jetson) & Veger T100 100~W USB-C PD bank + PD trigger cable
(15~V, 5.5$\times$2.5~mm barrel) & Independent 9--20~V supply for Orin mainboard \\
Power (motors) & 18650 cells (3S, $\sim$12.6~V) via JETANK expansion board &
Motor/servo rail only \\
Sensors & INA219 current sensor & Battery voltage/current monitoring \\
Networking & Intel Wireless 8265/8275 (PCIe) & WiFi via \texttt{wlP1p1s0}
(requires DKMS backport on JetPack~6.x) \\
\bottomrule
\end{tabular}
\end{table}

\section{Software Stack}
\label{sec:software-stack}

\begin{table}[htbp]
\centering
\caption{Software stack.}
\label{tab:software}
\begin{tabular}{@{}ll@{}}
\toprule
\textbf{Component} & \textbf{Version / Detail} \\
\midrule
OS & JetPack~6.2 (L4T R36.4.4, kernel 5.15.148-tegra) \\
GPU compute & CUDA~12.6, cuDNN~9.3.0.75, TensorRT~10.3.0.30 \cite{nvidia_tensorrt} \\
ML framework & PyTorch via Docker containers (moatazsawi/jetbot-orin) \cite{moatazsawi_jetbot_orin} \\
Computer vision & OpenCV~4.8.0, GStreamer with \texttt{nvarguscamerasrc} \\
Motor control & PCA9685 driver (\texttt{smbus2}), TB6612FNG via custom \texttt{Robot} class \\
Servo control & Feetech SCS SDK (\texttt{SCSCtrl}), 1~Mbps on \texttt{/dev/ttyTHS1} \\
Display & \texttt{luma.oled} (SSD1306, I2C bus~7), systemd service \\
Container runtime & Docker~28.2.2 \\
\bottomrule
\end{tabular}
\end{table}

%==============================================================================
\chapter[Methodology]{Methodology and Project Timeline}
\label{chap:methodology}
%==============================================================================

The project proceeded as a series of bring-up sessions, each resolving one or more
blocking issues before the next layer of the stack could be attempted. This chapter
summarises the chronological path; full session notes are retained in the project
archive.

\section{March --- Platform Misidentification and OS Selection}
\label{sec:march}

Early sessions assumed the hardware was a Jetson Nano 4~GB and attempted to flash
JetPack~4.x/5.x images accordingly. All such images failed to boot. The team
subsequently identified the actual hardware as a Jetson Orin Nano Super
(module~P3767, carrier board~P3768) and determined that JetPack~6.0+ is required
for this SoC generation. The earlier JetPack lines target the Tegra~X1 architecture
used in the original Nano, not the Orin's Tegra234/Ampere architecture. The JetPack~6.2
SD card image was the first to boot successfully.

\section{April --- WiFi, Repository Research, and Power Discovery}
\label{sec:april}

With SSH access established over Ethernet, the team diagnosed the WiFi blocker:
the onboard Intel Wireless~8265/8275 card was physically present on the PCIe bus
but produced no network interface because JetPack~6.x ships a stripped-down kernel
without the \texttt{iwlwifi} module. The team installed the \texttt{backport-iwlwifi-dkms}
package to resolve this. During hardware reassembly, the team discovered that the
JETANK expansion board cannot power the Jetson Orin Nano mainboard at all, which
fundamentally reshaped the project's power architecture (Section~\ref{sec:power-architecture}).

\section{Late April --- WiFi Resolution and Python Environment}
\label{sec:late-april}

The DKMS backport was explicitly installed against the running kernel, bringing up
the \texttt{wlP1p1s0} interface permanently. In parallel, the team resolved a broken
apt-managed pip installation by bootstrapping pip directly from the PyPA install
script, after a virtual-environment workaround proved insufficient for system-wide
package management.

\section{May --- Desktop Environment and System Stabilisation}
\label{sec:may}

A silent GDM3 failure (the \texttt{gdm3} binary itself was missing from the JetPack
image, masked by a systemd \texttt{RemainAfterExit} flag) was diagnosed and fixed by
reinstalling the package. With WiFi, pip, and the desktop environment all stable,
the team turned to full hardware assembly: camera, OLED, servos, and finally drive
motors.

\section{Hardware Bring-Up --- Camera, OLED, and Servos}
\label{sec:hardware-bringup}

The CSI camera, OLED status display, and serial bus servos were each brought online
during this phase, each requiring a JetPack~6.2-specific fix (device-tree overlay,
I2C bus renumbering, and calibration-value correction respectively). These are
detailed fully in Chapter~\ref{chap:task-c}.

\section{Drive Motor Investigation (Ongoing at Time of Reporting)}
\label{sec:motor-investigation}

The final and most significant blocker was the drive motor system: all motion commands
executed without error but produced no physical movement. Two confirmed software bugs
were fixed (Section~\ref{sec:drive-motor-bringup}), but the root cause was not fully
resolved before the reporting deadline, which in turn blocked implementation of
Tasks~B and~A.

%==============================================================================
\chapter{Results: Task C --- Hardware Assembly and Integration}
\label{chap:task-c}
%==============================================================================

Task~C covered chassis assembly, Jetson platform integration, and validation of
every onboard subsystem (camera, display, servos, motors) plus a stable remote
development workflow. Ten distinct system-level issues were identified and resolved
over the course of the semester; they are summarised in Table~\ref{tab:issues-summary}
and detailed in the following sections.

\begin{table}[htbp]
\centering
\caption{Summary of ten resolved issues in Task~C.}
\label{tab:issues-summary}
\small
\begin{tabular}{@{}p{0.6in}p{1.3in}p{1.6in}p{1.6in}@{}}
\toprule
\textbf{Category} & \textbf{Problem} & \textbf{Root Cause} & \textbf{Fix} \\
\midrule
OS & Image would not boot & JetPack~4.x/5.x images target Tegra~X1 (Nano),
not Tegra234/Ampere (Orin) & Flashed JetPack~6.2 SD card image \\
Networking & No WiFi interface present & \texttt{iwlwifi} kernel module stripped
from JetPack~6.x kernel & Installed \texttt{backport-iwlwifi-dkms}; explicit
\texttt{dkms install} step against running kernel \\
Desktop & No GUI on boot
(GDM3 ``active (exited)'') & \texttt{/usr/sbin/gdm3} binary missing from JetPack
image; \texttt{RemainAfterExit} masked failure & \texttt{apt install -\/\/reinstall gdm3} \\
Python & \texttt{pip} broken via apt & Version conflict between
\texttt{python3-setuptools} and \texttt{python3-pkg-resources} & Bootstrapped pip
via PyPA install script (bypassing apt) \\
Camera & No \texttt{/dev/video0} device & IMX219 driver activated via device-tree
overlay on JetPack~6.2, not a loadable module & Added \texttt{imx219-A.dtbo}
overlay to \texttt{extlinux.conf} \\
Display & OLED stayed blank & Hardcoded I2C bus~1, hardcoded \texttt{wlan0}
interface, unavailable libraries & Rewrote \texttt{stats.py} using
\texttt{luma.oled}, I2C bus~7, \texttt{wlP1p1s0}; deployed as systemd service \\
Servos & Aggressive spin, cable tangling & Centre positions drifted from factory
default of 512 after mounting & Read back centres via serial scan; updated
\texttt{servoInit} in installed package and source \\
Servos & Calibration script crashed silently & \texttt{servoInt.py} used
\texttt{time.sleep()} without \texttt{import time} & Added the missing import \\
Motors & PCA9685 channel forced HIGH & \texttt{set\_pwm(channel, 0, 0)} sets
ON=OFF=0; ON has priority (brake mode) & Changed to
\texttt{set\_pwm(channel, 0, 4096)} for correct FULL\_OFF \\
Motors & Wrong PCA9685 channel assignment & \texttt{motor.py} used Adafruit
MotorHAT layout instead of Waveshare JETANK layout & Remapped to Waveshare
standard (PWMA=0/AIN1=1/AIN2=2, PWMB=5/BIN1=3/BIN2=4) \\
\bottomrule
\end{tabular}
\end{table}

\section{Platform Identification and OS Flashing}
\label{sec:platform-id}

The single highest-leverage finding of the semester was platform identification:
the Jetson Orin Nano Super requires JetPack~6.0 or later \cite{nvidia_jetpack}. Several legacy SD card
images were tried and all failed to boot:
\texttt{jetbot-043\_nano-4gb-jp45},
\texttt{jetson-nano-jp461-sd-card-image},
\texttt{jp60dp-orin-nano-sd-card-image}, and
\texttt{JP512-orin-nano-sd-card-image};
only the JetPack~6.2 SD card image succeeded.
Flashing required Balena Etcher rather than Rufus, since Rufus cannot write raw
Linux partition images, and a Diskpart clean of the SD card on Windows before
flashing.

\section{Desktop Environment (GDM3)}
\label{sec:gdm3}

A particularly deceptive failure: \texttt{systemctl status} reported the display
manager as ``active (exited)''; normally a red flag would be more obvious. The
root cause was that the \texttt{gdm3} binary itself had never been installed in the
JetPack image; the init script's executable check silently exits when the binary
is missing, and a \texttt{RemainAfterExit=yes} flag in the systemd unit made the
service appear to have run successfully. Reinstalling the package resolved it.

\section{CSI Camera (IMX219-160)}
\label{sec:csi-camera}

Unlike most kernel drivers, the IMX219 sensor driver on JetPack~6.2 is enabled via
a device-tree overlay rather than a loadable kernel module. Without an
\texttt{OVERLAYS} line referencing the correct \texttt{.dtbo} file in
\texttt{/boot/extlinux/extlinux.conf}, the video input hardware initialises but
finds no configured sensor, and no \texttt{/dev/video*} device is created. Adding
the \texttt{imx219-A} overlay (for CAM0) resolved this; the \texttt{-C} variant
is reserved for CAM1 if the ribbon cable is ever moved.

\section{OLED Display}
\label{sec:oled}

The original JetBot \texttt{stats.py} script depended on three assumptions that do
not hold on this hardware: I2C bus~1 (the JETANK expansion board uses bus~7),
the \texttt{wlan0} interface name, and the \texttt{qwiic} / \texttt{Adafruit\_SSD1306}
libraries (not installed on JetPack~6.2 and effectively unmaintained). The script
was rewritten using \texttt{luma.oled} and \texttt{psutil}, and deployed as a
systemd service so the display comes up automatically on every boot.

\section{Power Architecture Mismatch}
\label{sec:power-architecture}

The JETANK expansion board was designed for the original Jetson Nano, which accepts
5~V via its GPIO header. The Jetson Orin Nano Super requires 9--20~V delivered
through a DC barrel jack and cannot be powered through GPIO at all. This meant the
two boards had to be powered through entirely separate paths: a Veger T100 USB-C
Power Delivery bank (set to 15~V via a PD trigger cable into a 5.5$\times$2.5~mm
barrel plug) for the Jetson, and the existing 18650 battery pack for the motor/servo
rail through the expansion board. The two rails share only a common ground through
the GPIO header, which is sufficient for I2C and serial communication.

One behavioural quirk: GPIO pins~2 and~4 are always live at 5~V whenever the Jetson
itself is powered, regardless of the expansion board's switch state, because the
OLED, INA219, and PCA9685 are all powered from the Jetson side, not the battery
side.

This power solution was identified and specified in detail but had not been procured
and tested end-to-end before the reporting deadline. This is recorded as Blocker~1
in Chapter~\ref{chap:discussion}.

\section{Servo Calibration}
\label{sec:servo-calibration}

Feetech SCS15-AP serial bus servos (IDs~1--5, daisy-chained on
\texttt{/dev/ttyTHS1}) initially spun aggressively to incorrect positions
on first command, tangling the base-rotation cable around the chassis.
The cause was that the servos' true centre positions had drifted from the factory
default of 512 after physical mounting. Reading back the actual centre positions
via a serial bus scan and updating \texttt{servoInit} to
\texttt{[None, 1021, 512, 512, 512, 811]} in both the installed package
and the source repository resolved the issue.

A second, unrelated bug (a missing \texttt{import time} in the calibration
script) was also fixed. Servo~1 (base rotation) has no physical end-stop and
is now constrained in software to $\pm$80$^\circ$ from centre to prevent future
cable tangling.

\section{Drive Motor Bring-Up}
\label{sec:drive-motor-bringup}

Two confirmed bugs were found and fixed in the \texttt{jetbot\_waveshare} Python
package:

\begin{enumerate}
\item \textbf{PCA9685 brake-mode bug:} the driver used
\texttt{set\_pwm(channel, 0, 0)} to represent a fully-off PWM value, but per the
PCA9685 datasheet, when ON=OFF=0 the ON bit takes priority and the pin is forced
permanently HIGH, keeping the TB6612FNG's direction pin high regardless of the
commanded duty cycle. This was corrected to
\texttt{set\_pwm(channel, 0, 4096)}, the documented FULL\_OFF encoding.
\item \textbf{Incorrect channel map:} the code used the Adafruit MotorHAT layout
(channels~2--4 for motor~1, 5--7 for motor~2) inherited from the original JetBot
codebase, instead of the Waveshare JETANK's actual layout (PWMA=0/AIN1=1/AIN2=2
for motor~1, PWMB=5/BIN1=3/BIN2=4 for motor~2).
\end{enumerate}

Neither fix produced motor movement. This is recorded as Blocker~2 and discussed
in full in Chapter~\ref{chap:discussion}.

%==============================================================================
\chapter{Results: Task B --- Camera Fusion for Perception}
\label{chap:task-b}
%==============================================================================

Task~B's objective was to produce a single fused perception state per frame,
combining a lane-following baseline (from the \texttt{road\_following} reference
notebooks) with an object-detection baseline (YOLO, from the
\texttt{object\_following} reference notebooks). The fusion design was completed
at the architectural level; implementation and on-hardware testing were blocked
by the unresolved power and motor issues described in Chapters~\ref{chap:task-c}
and~\ref{chap:discussion}.

\section{Fused State Schema}
\label{sec:fused-schema}

The designed per-frame fusion output combines a continuous lane-offset/confidence
signal with a discrete obstacle detection, plus a derived risk level and an action
hint intended for direct consumption by the Task~A controller:

\begin{lstlisting}[caption=Fused perception state schema, label=lst:fused-state,language={}]
{
  "lane_offset": 0.10,
  "lane_conf": 0.84,
  "obstacle": {"class": "person", "conf": 0.92, "distance_m": 1.7},
  "risk_level": "high",
  "action_hint": "stop"
}
\end{lstlisting}

\section{Planned Robustness Evaluation}
\label{sec:robustness-evaluation}

The evaluation plan called for testing the fused pipeline under five degraded
conditions, with frame-by-frame logging of FPS, confidence stability, and
qualitative failure cases for each:

\begin{itemize}
\item Lighting change
\item Shadows and glare
\item Motion blur
\item Partial occlusion of the lane or the obstacle
\item Different camera angles and resolutions
\end{itemize}

\section{Status}
\label{sec:task-b-status}

Not started on hardware. The lane-following and object-detection reference
notebooks (\texttt{road\_following}, \texttt{object\_following}) were identified
and reviewed as the intended baselines, and the Docker images for the ML and
display containers were pulled and verified runnable, but live fusion testing
requires a driveable robot, which was not achieved this semester.

%==============================================================================
\chapter{Results: Task A --- Camera Fusion for Driving}
\label{chap:task-a}
%==============================================================================

Task~A's objective was to convert the Task~B fused perception state into safe,
real-time driving commands. As with Task~B, the controller logic was fully designed
but not implemented or tested on hardware.

\section{Controller Policy (Designed)}
\label{sec:controller-policy}

\begin{itemize}
\item Lane centred and clear $\rightarrow$ move forward
\item Lane offset detected $\rightarrow$ apply steering correction proportional
to the offset
\item High risk level or low perception confidence $\rightarrow$ slow down or stop
\end{itemize}

\section{Safety Gates (Designed)}
\label{sec:safety-gates}

\begin{itemize}
\item \textbf{Confidence-threshold stop:} if lane or obstacle confidence falls
below a set threshold, halt rather than act on an uncertain perception state
\item \textbf{Obstacle-distance-threshold stop:} halt if a detected obstacle is
closer than a safe following distance
\item \textbf{Manual emergency-stop path:} an independent, software-prioritised
stop command that overrides the controller output
\end{itemize}

\section{Closed-Loop Evaluation Plan (Designed)}
\label{sec:evaluation-plan}

Three scenarios were planned for quantitative and qualitative evaluation once the
robot could drive:

\begin{enumerate}
\item Normal lane following
\item An obstacle appearing in the robot's path
\item An ambiguous, low-confidence perception case
\end{enumerate}

For each scenario, the plan calls for tracking intervention count, command
stability/jitter, and safe-stop trigger reliability.

\section{Status}
\label{sec:task-a-status}

Not started on hardware. This task is fully gated on Task~B perception output and
a working drive-motor system (Chapter~\ref{chap:discussion}, Blocker~2).

%==============================================================================
\chapter{Discussion of Critical Blockers}
\label{chap:discussion}
%==============================================================================

\section{Blocker 1 --- Power Architecture}
\label{sec:blocker1}

The JETANK expansion board cannot power the Jetson Orin Nano under any configuration;
this is a fixed property of the hardware, not a bug. A specific solution (a
Veger T100 100~W USB-C PD power bank with a 15~V trigger cable into the Orin's
barrel jack) was identified, with the constraint that the connection must be
USB-C to USB-C (a USB-A source bypasses PD negotiation entirely and the bank falls
back to 5~V, which will not boot the Jetson). This was not procured and tested in
time to support full untethered operation.

\section{Blocker 2 --- Drive Motors Unresponsive}
\label{sec:blocker2}

This is the project's most significant open technical problem. After fixing the two
confirmed software bugs (Section~\ref{sec:drive-motor-bringup}), the team ruled out
the following as causes through direct testing:

\begin{itemize}
\item PCA9685 not present on the I2C bus: ruled out; confirmed at address
\texttt{0x60} on bus~7 via \texttt{i2cdetect}
\item Insufficient battery voltage: ruled out; INA219 reported 11.85~V
(79\% charge) during testing
\item Incorrect motor wiring to the expansion board terminals: ruled out;
connections were physically verified
\item Wrong Python import path: ruled out; \texttt{from jetbot import Robot}
was confirmed to resolve to the corrected \texttt{jetbot\_waveshare} package
\item I2C communication failures: ruled out; no I2C errors were thrown
on any \texttt{set\_pwm} call
\end{itemize}

An exhaustive brute-force scan of plausible PCA9685 channel-to-motor mappings,
including STBY-candidate channels~8,~9,~12,~13,~14, and~15, also produced no
movement on any combination.

What remains unverified:
\begin{itemize}
\item Whether PCA9685 register writes are actually sticking (no register read-back
has been performed)
\item The exact PCA9685-to-TB6612FNG pin wiring on the JETANK expansion board
(no official schematic is publicly available)
\item Whether the TB6612FNG's STBY pin is hardwired high or requires explicit
activation from an as-yet-untried channel
\item Whether the TB6612FNG's VM pin is actually receiving the 12~V motor-supply
rail through the expansion board's traces
\end{itemize}

The working hypothesis is that the fault sits below the software layer the team
controls, most likely a wiring or enable-pin issue on the expansion board itself,
only resolvable through direct multimeter probing or a manufacturer-provided
schematic.

\section{Downstream Effect on Tasks B and A}
\label{sec:downstream-effect}

Because Tasks~B and~A both assume a robot capable of driving, neither could move
past the design stage this semester. The perception-fusion schema
(Section~\ref{sec:fused-schema}) and the safety-gated controller policy
(Sections~\ref{sec:controller-policy}--\ref{sec:safety-gates}) represent the team's
complete intended design and are ready to be implemented as soon as Blockers~1
and~2 are resolved.

%==============================================================================
\chapter{Conclusion}
\label{chap:conclusion}
%==============================================================================

This project set out to build an autonomous, camera-fusion-driven JetBot on a
Jetson Orin Nano Super. The headline technical result is not the driving demo
originally scoped, but a thorough, hard-won account of why an Orin-generation Jetson
is not a drop-in replacement for the Jetson Nano the JETANK kit and most public
JetBot tutorials were designed around, and a documented, ten-issue bring-up path
that resolves every layer of that mismatch except the motor driver hardware itself.

Ten distinct system-level issues were diagnosed and fixed across OS imaging, kernel
driver support, desktop environment, Python packaging, CSI camera enablement, OLED
display, and serial servo calibration. Two blockers remain open: an unprocured
power-bank solution for the Jetson mainboard, and an unresolved drive-motor fault
that has been narrowed to the hardware layer through a careful process of elimination.
The perception-fusion and driving-control logic for Tasks~B and~A were fully designed
and are implementation-ready.

Although the robot did not drive by the end of the semester, the project produced
a complete and transferable diagnostic record: every fix is traceable to a specific
root cause, and every open item has a specific, actionable next diagnostic step
(Chapter~\ref{chap:recommendations}). That record is, itself, the project's primary
deliverable.

%==============================================================================
\chapter[Recommendations]{Recommendations and Future Work}
\label{chap:recommendations}
%==============================================================================

\section{To Close Out the Drive Motor Issue}
\label{sec:close-motor}

\begin{enumerate}
\item Read back PCA9685 LED0--15 registers via \texttt{i2cget} after each write
to confirm values are actually sticking in hardware, ruling in or out a silent
I2C write failure.
\item Probe the TB6612FNG directly with a multimeter while motion commands are
issued: VM (motor supply voltage), STBY (should be above 2~V for the H-bridge to
be enabled), and IN1/IN2 (should toggle between 0~V and 3.3~V with commanded
direction).
\item Expand the channel brute-force search to include PCA9685 channels~8--15 in
every combination tried, not only the previously assumed STBY candidates.
\item Temporarily wire one motor directly to the battery pack to rule out a damaged
motor before continuing to debug the driver chain.
\item Contact Waveshare support directly for the JETANK expansion board schematic,
or at minimum the PCA9685-to-TB6612FNG pin mapping.
\end{enumerate}

\section{To Close Out the Power Architecture}
\label{sec:close-power}

\begin{enumerate}
\item Procure the Veger T100 and the 15~V PD trigger cable (5.5$\times$2.5~mm barrel)
and test the full wireless-plus-battery configuration end-to-end.
\item Evaluate the Waveshare UPS Module (C), built specifically for the Jetson Orin
family around 21700 cells, as a cleaner single-package alternative.
\item Mount the chosen power source low and centred on the chassis for balance once
finalised.
\end{enumerate}

\section{For Future Orin-Generation JetBot / JETANK Projects}
\label{sec:future-projects}

\begin{enumerate}
\item Use JetPack~6.2 (or later) as the baseline image from day one; earlier
JetPack lines are incompatible with Orin hardware at the SoC level.
\item Add the IMX219 device-tree overlay and install \texttt{backport-iwlwifi-dkms}
as standard first-boot steps immediately after flashing.
\item Replace every \texttt{wlan0} reference in inherited scripts with
\texttt{wlP1p1s0}; JetPack~6.x uses predictable network interface naming.
\item Confirm the I2C bus number for the expansion board before writing any
peripheral code; it is bus~7 on this hardware, not the bus~1 most tutorials
assume.
\item Bootstrap pip directly from the PyPA install script rather than relying on
apt-managed pip, which is broken out of the box on JetPack~6.x.
\item Plan for an NVMe SSD upgrade early; the 28~GiB microSD card fills quickly
once Docker images and notebooks are added (83\% used before any ML datasets were
stored).
\end{enumerate}

\section{Once Tasks B and A Are Unblocked}
\label{sec:when-unblocked}

\begin{enumerate}
\item Implement and test the Task~B fusion pipeline against the planned robustness
scenarios (lighting, shadows/glare, blur, occlusion, camera angle/resolution
changes) before integrating with Task~A.
\item Implement the Task~A controller exactly as designed in
Chapter~\ref{chap:task-a}, then run the three planned closed-loop scenarios and
report intervention count, command stability, and safe-stop reliability as the
project's quantitative result set.
\end{enumerate}

%==============================================================================
\chapter{Self-Reflection Note}
\label{chap:self-reflection}
%==============================================================================

\section*{Individual Contribution Breakdown}

\begin{table}[htbp]
\centering
\caption{Individual contribution percentages.}
\label{tab:contributions}
\large
\begin{tabular}{@{}lc@{}}
\toprule
\textbf{Team Member} & \textbf{Contribution (\%)} \\
\midrule
Adam Ryan   & [\%] \\
Aidan Ryan  & [\%] \\
Denis Vasilev & [\%] \\
Nora Regina Cerveny & [\%] \\
\midrule
Total & 100\% \\
\bottomrule
\end{tabular}
\end{table}

\newpage
\section*{Adam Ryan}
[Write your self-reflection here --- approximately 8--10 lines. What did you learn?
What challenges did you face? How did you contribute to the team? What would you
do differently?]

\section*{Aidan Ryan}
[Write your self-reflection here --- approximately 8--10 lines. What did you learn?
What challenges did you face? How did you contribute to the team? What would you
do differently?]

\section*{Denis Vasilev}
[Write your self-reflection here --- approximately 8--10 lines. What did you learn?
What challenges did you face? How did you contribute to the team? What would you
do differently?]

\section*{Nora Regina Cerveny}
[Write your self-reflection here --- approximately 8--10 lines. What did you learn?
What challenges did you face? How did you contribute to the team? What would you
do differently?]

%%   BACK MATTER  %%%%%%%%%%%%%%%%%%%%%%%%%%%%%%%%%%%%%%%%%%%%%%%%%%%%%%%%%%%%%%
\backmatter%

\addtoToC{Bibliography}
\bibliographystyle{IEEEtranS}
\typeout{}
\bibliography{references}

%%   APPENDICES  %%%%%%%%%%%%%%%%%%%%%%%%%%%%%%%%%%%%%%%%%%%%%%%%%%%%%%%%%%%%%%%
\begin{appendices}

\chapter{Full Resolved-Issue Log}
\label{app:issue-log}

This appendix expands Table~\ref{tab:issues-summary} with the blocking status of
each issue, for archival and grading reference.

\begin{longtable}{@{}cp{1.8in}p{1.8in}p{1.2in}@{}}
\caption{Complete issue log with blocking status.}
\label{tab:issue-log}\\
\toprule
\textbf{\#} & \textbf{Issue} & \textbf{Resolution} & \textbf{Blocker for Driving?} \\
\midrule
\endfirsthead
\multicolumn{4}{c}{\tablename\ \thetable{} --- continued from previous page} \\
\toprule
\textbf{\#} & \textbf{Issue} & \textbf{Resolution} & \textbf{Blocker for Driving?} \\
\midrule
\endhead
\bottomrule
\endfoot
1 & Wrong OS for Orin hardware (JP4.x/5.x images) & JetPack~6.2 SD card image
& Yes --- resolved \\
2 & No WiFi (\texttt{iwlwifi} missing from kernel) & \texttt{backport-iwlwifi-dkms}
& Yes --- resolved \\
3 & GDM3 binary missing &
\texttt{apt install --reinstall gdm3} & No \\
4 & \texttt{pip} broken via apt &
Bootstrapped pip via PyPA install script & No \\
5 & Camera not detected (no \texttt{/dev/video0}) &
IMX219 device-tree overlay in \texttt{extlinux.conf} & No \\
6 & OLED: wrong I2C bus, interface name, missing libs &
Rewrote with \texttt{luma.oled}, bus~7, \texttt{wlP1p1s0} & No \\
7 & Servo centres drifted from factory default &
Updated \texttt{servoInit} in both \texttt{TTLServo.py} copies & No \\
8 & \texttt{servoInt.py} missing \texttt{import time} &
Added the missing import & No \\
9 & PCA9685 \texttt{set\_pwm(0,0)} forces pin HIGH (brake mode) &
Changed to \texttt{set\_pwm(0, 4096)} for FULL\_OFF &
Yes --- critical (fix applied, motors still not moving) \\
10 & Motor channel mapping wrong for Waveshare hardware &
Updated to Waveshare PCA9685 channel layout &
Yes --- critical (fix applied, motors still not moving) \\
\end{longtable}

\end{appendices}
\end{document}
